% GENERATED FILE. Edit canonical lessons, then run npm run generate:books.
% canonical-source-hash: fnv1a64:7c59d015b09ca9c8
% canonical-lessons: ZH-C441-jiaoqu, ZH-C441-yewai, ZH-C441-zhangpeng, ZH-C441-shenghuo, ZH-C441-tianran, ZH-R441-review-436-441

\chapter{Out in the Open}
\label{ch:zh-out-in-the-open}
\hlchaptermodality{441}

\begin{chapteropening}
\textbf{By the end of this chapter:} \emph{I can talk about the outskirts, the open country, a tent, lighting a fire and a lake made by nature.}

Complete the last lesson of chapter 441: I can talk about the outskirts, the open country, a tent, lighting a fire and a lake made by nature.
\end{chapteropening}

\section[jiāoqū]{jiāoqū — suburbs; outskirts}

\label{lesson:ZH-C441-jiaoqu}



\begin{quote}
Before the new one: say the Mandarin for a villager, then the Mandarin for to lodge with a local family.
\end{quote}

\subsection*{You'll want to know: jiāoqū}
\emph{jiāoqū} — \textquotedblleft{}suburbs; outskirts\textquotedblright{}.

\begin{quote}
\emph{wǒmen zhù zài jiāoqū zuò dìtiě qù shìzhōngxīn} — We live in the suburbs and take the underground into the centre.
\end{quote}

Its characters come later; for now it is a word for the ear.

\begin{cousinweb}
\emph{jiāoqū} is built from \emph{jiāo} 'the land beyond the walls' and \zh{区} 'an area', both inherited from Middle Chinese, \emph{kaew} and \emph{khju} (Guangyun, in Baxter's transcription), which glosses the first as the land outside a walled town.
\end{cousinweb}

\begin{grammarlens}[title={What you've built}]
One more word for getting around.
\end{grammarlens}

\subsection*{Your turn}
\begin{itemize}
\raggedright
  \item \emph{Say it:} \emph{jiāoqū}
  \item \emph{Say it:} \emph{jiāoqū}, once more
  \item \emph{Say it:} \emph{wǒmen zhù zài jiāoqū zuò dìtiě qù shìzhōngxīn} — We live in the suburbs and take the underground into the centre.
  \item \emph{Recall:} say the Mandarin for a villager, then the Mandarin for to lodge with a local family, then say \emph{jiāoqū} again
\end{itemize}

\begin{tcolorbox}[breakable,colback=teal!4,colframe=teal!35!black,title={Before you move on}]
What does jiāoqū mean? (\textquotedblleft{}Suburbs; outskirts\textquotedblright{}.) Say it once more.
\end{tcolorbox}
\section[yěwài]{yěwài — the open country; the wild}

\label{lesson:ZH-C441-yewai}



\begin{quote}
Before the new one: say the Mandarin for to lodge with a local family, then the Mandarin for the outskirts of a city.
\end{quote}

\subsection*{You'll want to know: yěwài}
\emph{yěwài} — \textquotedblleft{}the open country; the wild\textquotedblright{}.

\begin{quote}
\emph{zhōumò wǒmen zuò chē qù yěwài} — At the weekend we're driving out into the open country.
\end{quote}

Its characters come later; for now it is a word for the ear.

\begin{cousinweb}
\emph{yěwài} is built from \emph{yě} 'open fields' and \zh{外} 'outside', both inherited from Middle Chinese, \emph{yaeX} and \emph{ngwajH} (Guangyun, in Baxter's transcription), which glosses the first, after an older dictionary, as the land beyond the town.
\end{cousinweb}

\begin{grammarlens}[title={What you've built}]
One more word for getting around.
\end{grammarlens}

\subsection*{Your turn}
\begin{itemize}
\raggedright
  \item \emph{Say it:} \emph{yěwài}
  \item \emph{Say it:} \emph{yěwài}, once more
  \item \emph{Say it:} \emph{zhōumò wǒmen zuò chē qù yěwài} — At the weekend we're driving out into the open country.
  \item \emph{Recall:} say the Mandarin for to lodge with a local family, then the Mandarin for the outskirts of a city, then say \emph{yěwài} again
\end{itemize}

\begin{tcolorbox}[breakable,colback=teal!4,colframe=teal!35!black,title={Before you move on}]
What does yěwài mean? (\textquotedblleft{}The open country; the wild\textquotedblright{}.) Say it once more.
\end{tcolorbox}
\section[zhàngpeng]{zhàngpeng — tent}

\label{lesson:ZH-C441-zhangpeng}



\begin{quote}
Before the new one: say the Mandarin for the outskirts of a city, then the Mandarin for out in the wild.
\end{quote}

\subsection*{You'll want to know: zhàngpeng}
\emph{zhàngpeng} — \textquotedblleft{}tent\textquotedblright{}.

\begin{quote}
\emph{wǎnshang wǒmen zài zhàngpeng lǐ shuìjiào} — At night we sleep in the tent.
\end{quote}

Its characters come later; for now it is a word for the ear.

\begin{cousinweb}
\emph{zhàngpeng} is built from \emph{zhàng} 'a curtain' and \emph{péng} 'an awning', both inherited from Middle Chinese, \emph{trjangH} and \emph{buwng} (Guangyun, in Baxter's transcription), which glosses the second as a woven cover over a boat; in this word it is said light.
\end{cousinweb}

\begin{grammarlens}[title={What you've built}]
One more word for getting around.
\end{grammarlens}

\subsection*{Your turn}
\begin{itemize}
\raggedright
  \item \emph{Say it:} \emph{zhàngpeng}
  \item \emph{Say it:} \emph{zhàngpeng}, once more
  \item \emph{Say it:} \emph{wǎnshang wǒmen zài zhàngpeng lǐ shuìjiào} — At night we sleep in the tent.
  \item \emph{Recall:} say the Mandarin for the outskirts of a city, then the Mandarin for out in the wild, then say \emph{zhàngpeng} again
\end{itemize}

\begin{tcolorbox}[breakable,colback=teal!4,colframe=teal!35!black,title={Before you move on}]
What does zhàngpeng mean? (\textquotedblleft{}Tent\textquotedblright{}.) Say it once more.
\end{tcolorbox}
\section[\texorpdfstring{\zh{生火} (shēnghuǒ)}{生火 (shēnghuǒ)}]{\zh{生火} (shēnghuǒ) — to light a fire}

\label{lesson:ZH-C441-shenghuo}



\begin{quote}
Before the new one: say the Mandarin for out in the wild, then the Mandarin for a tent.
\end{quote}

\subsection*{You'll want to know: \zh{生火}}
\textbf{\zh{生火}} — \emph{shēnghuǒ} — \textquotedblleft{}to light a fire\textquotedblright{}.

\begin{quote}
\textbf{\zh{天冷了我们在村子外边生火}} \emph{tiān lěng le wǒmen zài cūnzi wàibian shēnghuǒ} — It's got cold, so we're lighting a fire outside the village.
\end{quote}

\begin{cousinweb}
\textbf{\zh{生火}} (\emph{shēnghuǒ}) is built from \zh{生} 'to bring to life' and \zh{火} 'fire', both inherited from Middle Chinese, \emph{srjaeng} and \emph{xwaX} (Guangyun, in Baxter's transcription), which glosses the first 'to grow'.
\end{cousinweb}

\begin{grammarlens}[title={What you've built}]
One more word for getting around.
\end{grammarlens}

\subsection*{Your turn}
\begin{itemize}
\raggedright
  \item \emph{Say it:} \emph{shēnghuǒ}
  \item \emph{Say it:} \emph{shēnghuǒ}, once more
  \item \emph{Say it:} \emph{tiān lěng le wǒmen zài cūnzi wàibian shēnghuǒ} — It's got cold, so we're lighting a fire outside the village.
  \item \emph{Recall:} say the Mandarin for out in the wild, then the Mandarin for a tent, then say \emph{shēnghuǒ} again
\end{itemize}

\begin{tcolorbox}[breakable,colback=teal!4,colframe=teal!35!black,title={Before you move on}]
What does \zh{生火} mean? (\textquotedblleft{}To light a fire\textquotedblright{}.) Say it once more.
\end{tcolorbox}
\section[\texorpdfstring{\zh{天然} (tiānrán)}{天然 (tiānrán)}]{\zh{天然} (tiānrán) — natural; made by nature}

\label{lesson:ZH-C441-tianran}



\begin{quote}
Before the new one: say the Mandarin for a tent, then the Mandarin for to light a fire.
\end{quote}

\subsection*{You'll want to know: \zh{天然}}
\textbf{\zh{天然}} — \emph{tiānrán} — \textquotedblleft{}natural; made by nature\textquotedblright{}.

\begin{quote}
\emph{zhè shì yí ge tiānrán de hú} — This is a natural lake.
\end{quote}

\begin{cousinweb}
\textbf{\zh{天然}} (\emph{tiānrán}) is built from \zh{天} 'heaven' and \zh{然} 'so', both inherited from Middle Chinese, \emph{then} and \emph{nyen} (Guangyun, in Baxter's transcription), which glosses the second as a particle and 'thus': so by heaven, not by people. How the pair came to mean 'natural' is not settled.
\end{cousinweb}

\begin{grammarlens}[title={What you've built}]
One more word for getting around.
\end{grammarlens}

\subsection*{Your turn}
\begin{itemize}
\raggedright
  \item \emph{Say it:} \emph{tiānrán}
  \item \emph{Say it:} \emph{tiānrán}, once more
  \item \emph{Say it:} \emph{zhè shì yí ge tiānrán de hú} — This is a natural lake.
  \item \emph{Recall:} say the Mandarin for a tent, then the Mandarin for to light a fire, then say \emph{tiānrán} again
\end{itemize}

\begin{tcolorbox}[breakable,colback=teal!4,colframe=teal!35!black,title={Before you move on}]
What does \zh{天然} mean? (\textquotedblleft{}Natural; made by nature\textquotedblright{}.) Say it once more.
\end{tcolorbox}
\section[(dialogue)]{Looking back: from North and South to Out in the Open}

\label{lesson:ZH-R441-review-436-441}



\begin{quote}
Say the Mandarin for the lie of the land and the south of the country.
\end{quote}

\subsection*{Your turn}
\begin{itemize}
\raggedright
  \item \emph{Say it:} the lie of the land, the south of the country, the north of the country, a regional dialect, a regional accent
  \item \emph{Say it:} a mountain area, flat open land, a desert, a high plateau, to emigrate
  \item \emph{Say it:} to open up land for use, the northern part, the western part, a part of a whole, all of it
  \item \emph{Say it:} a rice paddy, an orchard, a farm stay, a duck, to take one's mind off things
  \item \emph{Say it:} to go to market day, a mountain village, the village head, a villager, to lodge with a local family
  \item \emph{Say it:} the outskirts of a city, out in the wild, a tent, to light a fire, made by nature
  \item \emph{Read it:} \zh{方} \zh{部} \zh{村} — each one aloud, then a word that holds it
\end{itemize}

\begin{tcolorbox}[breakable,colback=teal!4,colframe=teal!35!black,title={Before you move on}]
The lie of the land? (\emph{dìxíng}.) The south of the country? (\textbf{\zh{南方}}, \emph{nánfāng}.) The north of the country? (\textbf{\zh{北方}}, \emph{běifāng}.) A regional dialect? (\emph{fāngyán}.) A regional accent? (\textbf{\zh{口音}}, \emph{kǒuyīn}.) A mountain area? (\textbf{\zh{山区}}, \emph{shānqū}.) Flat open land? (\emph{píngyuán}.) A desert? (\emph{shāmò}.) A high plateau? (\emph{gāoyuán}.) To emigrate? (\emph{yímín}.) To open up land for use? (\textbf{\zh{开发}}, \emph{kāifā}.) The northern part? (\textbf{\zh{北部}}, \emph{běibù}.) The western part? (\textbf{\zh{西部}}, \emph{xībù}.) A part of a whole? (\emph{bùfen}.) All of it? (\textbf{\zh{全部}}, \emph{quánbù}.) A rice paddy? (\emph{dàotián}.) An orchard? (\emph{guǒyuán}.) A farm stay? (\emph{nóngjiālè}.) A duck? (\emph{yāzi}.) To take one's mind off things? (\emph{sànxīn}.) To go to market day? (\emph{gǎnjí}.) A mountain village? (\textbf{\zh{山村}}, \emph{shāncūn}.) The village head? (\textbf{\zh{村长}}, \emph{cūnzhǎng}.) A villager? (\emph{cūnmín}.) To lodge with a local family? (\emph{jièsù}.) The outskirts of a city? (\emph{jiāoqū}.) Out in the wild? (\emph{yěwài}.) A tent? (\emph{zhàngpeng}.) To light a fire? (\textbf{\zh{生火}}, \emph{shēnghuǒ}.) Made by nature? (\textbf{\zh{天然}}, \emph{tiānrán}.)
\end{tcolorbox}
